\begin{lemma}\label{lem:2.2.11}
设 $(X,\sigma)$ 是克莱因曲面，$p:(\tilde X,\tilde x)\to(X,x)$ 是 $X$ 的万有覆盖，其中
$x$ 是 $X$ 中任一点。则
\begin{enumerate}[label=(\roman*), leftmargin=2.2em, itemsep=0.15em]
\item 对任意 $y\in p^{-1}(\{\sigma(x)\})$，存在唯一的反全纯映射
      $\tilde\sigma:\tilde X\to\tilde X$，使得 $\tilde\sigma(\tilde x)=y$
      且下图可交换：
\[
\begin{tikzcd}[column sep=large, row sep=large]
\tilde X \arrow[r, "\tilde\sigma"] \arrow[d, "p"'] & \tilde X \arrow[d, "p"]\\
X \arrow[r, "\sigma"'] & X
\end{tikzcd}
\]
\item 若 $x\in\Fix(\sigma)$ 且 $y:=\tilde x\in p^{-1}(\{x\})=p^{-1}(\{\sigma(x)\})$，则
      (i) 中定义的反全纯映射 $\tilde\sigma$ 满足 $\tilde\sigma^2=\Id_{\tilde X}$；
      换言之，它是 $\tilde X$ 上的一个实结构。
\end{enumerate}
\end{lemma}

\begin{proof}
本引理说的是：$\sigma$ 总能提升为 $\tilde X$ 的一个覆盖 $\sigma$ 的反全纯变换，但一般
不能提升为对合的反全纯变换。考虑（反全纯的）覆盖映射
$\sigma\circ p:(\tilde X,\tilde x)\to(X,\sigma(x))$。由于 $\tilde X$ 单连通，
$\sigma\circ p$ 是 $X$ 的拓扑万有覆盖；因此给定 $y\in p^{-1}(\{\sigma(x)\})$，存在唯一
的覆盖映射 $\tilde\sigma:(\tilde X,\tilde x)\to(\tilde X,y)$ 满足
$p\circ\tilde\sigma=\sigma\circ p$。由于 $(\tilde X,y)$ 单连通，覆盖映射 $\tilde\sigma$
必是单射，从而是微分同胚。此外，由于 $\sigma$ 反全纯而 $p$ 是全纯覆盖映射，关系
$p\circ\tilde\sigma=\sigma\circ p$ 蕴涵 $\tilde\sigma$ 反全纯。这就证明了 (i)。

若 $x\in\Fix(\sigma)$ 且 $y=\tilde x$，则 $\tilde\sigma^2:\tilde X\to\tilde X$ 覆盖
$\sigma^2=\Id_X$ 且把 $\tilde x$ 送到 $\tilde x$。由这种映射的唯一性，
$\tilde\sigma^2=\Id_{\tilde X}$，这就证明了 (ii)。
\end{proof}

注意上述 (ii) 表明：$X$ 中实点（若存在）的选取定义了 $\tilde X$ 上的一个实结构，它使
万有覆盖映射 $p:\tilde X\to X$ 成为克莱因曲面同态。$X$ 中不同的实点可能导致 $\tilde X$
上不同构的实结构（另见命题 2.2.13 之后的说明）。

引理 2.2.11 的第一个应用是实曲线基本群的一个性质。虽然它从我们对这个群的定义看是
显然的，但对后面的一切、特别是对 2.3 节将介绍的实曲线基本群的表示概念，都至关重要。

\begin{proposition}\label{prop:2.2.12}
设 $\alpha:\pi^R_1((X,\sigma);x)\to\{\Id_X;\sigma\}^{\mathrm{op}}\simeq\Z/2\Z$ 是
命题 2.2.8 中定义的群同态。则有短正合列，称为\emph{同伦正合列}，
\begin{equation}\label{eq:2.17}
1\longrightarrow\pi_1(X;x)\longrightarrow\pi^R_1((X,\sigma);x)
\xrightarrow{\ \alpha\ }\Z/2\Z\longrightarrow1,
\end{equation}
它特别说明 $\pi_1(X;x)$ 是 $\pi^R_1((X,\sigma);x)$ 的指数为 $2$ 的子群。
\end{proposition}

\begin{proof}
由命题 2.2.8，只需证明 $\alpha$ 是满射。这是引理 2.2.11(i) 的推论。
\end{proof}

引理 2.2.11 的下一个应用是证明：当基点 $x$ 是 $\sigma$ 的不动点时，群
$\pi^R_1((X,\sigma);x)$ 具有典范的半直积结构（另见习题 2.2.30 与 2.2.31）。

\begin{proposition}\label{prop:2.2.13}
设 $x\in\Fix(\sigma)$ 是 $X$ 的实点。由引理 2.2.11，存在 $\tilde X$ 上的实结构
$\tilde\sigma$ 把 $\tilde x$ 送到 $\tilde x$。这个 $\tilde\sigma$ 通过共轭给出
$\{\Id_{\tilde X};\tilde\sigma\}^{\mathrm{op}}\simeq\Z/2\Z$ 在
$\pi_1(X;x)=\Aut(\tilde X/X)^{\mathrm{op}}$ 上的左作用
\begin{equation}\label{eq:2.18}
\tilde\sigma\cdot f=\tilde\sigma f\tilde\sigma^{-1},
\end{equation}
并且 $\pi^R_1((X,\sigma);x)$ 同构于半直积
$\pi_1(X;x)\rtimes_{\tilde\sigma}\Z/2\Z$。
\end{proposition}

注意上述断言中隐含了：符号 $\cdot$ 表示 $\Z/2\Z$ 在 $\pi_1(X;x)$ 上的群作用，而
$\tilde\sigma f\tilde\sigma^{-1}$ 表示 $\tilde\sigma$、$f$ 与 $\tilde\sigma^{-1}$ 在
$\pi^R_1((X,\sigma);x)$ 中的乘积。从现在起，我们将使用这一记号，而不再使用此前所用的
$\tilde\sigma\curvearrowright f=\tilde\sigma\cdot f\cdot\tilde\sigma^{-1}
=\tilde\sigma^{-1}\circ f\circ\tilde\sigma$。

\begin{proof}[命题 2.2.13 的证明]
先注意：若 $f$ 覆盖 $\Id_X$，则 $\tilde\sigma f\tilde\sigma^{-1}$ 也覆盖 $\Id_X$；
若 $f$ 全纯，则它也是全纯的。这定义了 $\Z/2\Z$ 在 $\pi_1(X;x)$ 上的左作用，因为容易
验证对一切 $\varepsilon,\eta\in\{\Id_X,\sigma\}$ 有
$\widetilde{\varepsilon\eta}=\tilde\varepsilon\tilde\eta$。若
$g\in\pi^R_1((X,\sigma);x)$ 覆盖 $\sigma$，则 $g=(g\tilde\sigma^{-1})\tilde\sigma$，
其中 $g\tilde\sigma^{-1}$ 覆盖 $\Id_X$，故映射
\begin{equation}\label{eq:2.19}
\pi_1(X;x)\rtimes_{\tilde\sigma}\{\Id_{\tilde X};\tilde\sigma\}
\longrightarrow\pi^R_1((X,\sigma);x),
\qquad (f,\varepsilon)\longmapsto f\varepsilon
\end{equation}
是满射。又因对 $f\in\pi_1(X;x)$，$f\varepsilon$ 在 $X$ 上覆盖 $\sigma$ 当且仅当
$\varepsilon=\tilde\sigma$，故它也是单射。最后，乘积
$(f,\varepsilon)(g,\eta)=\bigl(f(\varepsilon\cdot g),\varepsilon\eta\bigr)
=(f\varepsilon g\varepsilon^{-1},\varepsilon\eta)$ 在
$\pi_1(X;x)\rtimes_{\tilde\sigma}\{\Id_{\tilde X};\tilde\sigma\}$ 中，经 (2.19) 被送到
$(f\varepsilon g\varepsilon^{-1})(\varepsilon\eta)=(f\varepsilon)(g\eta)$，故 (2.19)
确实是群同构。
\end{proof}

我们回忆：$\pi^R_1((X,\sigma);x)$ 上的半直积结构，与同伦正合列 (2.17) 的一个分裂
$\Z/2\Z\to\pi^R_1((X,\sigma);x)$ 是同一件事，而这也与
$\pi^R_1((X,\sigma);x)$ 中一个二阶反全纯元素，即 $\tilde X$ 上覆盖 $X$ 上 $\sigma$ 的
实结构，是同一件事。因此，引理 2.2.11 与命题 2.2.13 表明：$(X,\sigma)$ 的一个实点给出
同伦正合列
\[
1\longrightarrow\pi_1(X;x)\longrightarrow\pi^R_1((X,\sigma);x)
\xrightarrow{\ \alpha\ }\Z/2\Z\longrightarrow1
\]
的一个分裂。当然，不同的实点可能导致不同的半直积结构，但弄清这究竟如何发生并不是一件
初等的事，关于这方面的结果可参看 [13]。

引理 2.2.11 的最后一个应用，是用 $x$ 处的环路（以及当 $\sigma(x)\ne x$ 时从 $x$ 到
$\sigma(x)$ 的道路）来描述 $(X,\sigma)$ 的基本群。我们将看到，这一描述取决于 $x$
是否为 $(X,\sigma)$ 的实点而有所不同。这正是我们给出 $\pi^R_1((X,\sigma);x)$ 的一个
不涉及 $X$ 中道路的定义的原因，因为那给出了群 $\pi^R_1((X,\sigma);x)$ 的一个更内蕴的
描述，其中无须按 $x$ 是否为实点分情形讨论。当 $x\in\Fix(\sigma)$ 时，用 $x$ 处环路
对 $\pi^R_1((X,\sigma);x)$ 的描述其实是命题 2.2.13 的推论。

\begin{corollary}\label{cor:2.2.14}
把 $\pi_1(X;x)$ 想成 $x\in\Fix(\sigma)$ 处环路的同伦类群。则对任意
$\gamma\in\pi_1(X;x)$，(2.18) 中定义的 $\{\Id_X,\sigma\}^{\mathrm{op}}\simeq\Z/2\Z$
在 $\pi_1(X;x)$ 上的左作用由 $\sigma_*\gamma=\sigma^{-1}\circ\gamma$ 给出；于是
$\pi^R_1((X,\sigma);x)\simeq\pi_1(X;x)\rtimes_\sigma\Z/2\Z$，其中的半直积是就该作用
而言的。
\end{corollary}

\begin{proof}
由命题 2.2.13，只需证明 $\Psi(\sigma_*\gamma)=\tilde\sigma f\tilde\sigma^{-1}$，其中
$\Psi:\pi_1(X;x)\to\Aut(\tilde X/X)^{\mathrm{op}}$ 是命题 2.2.4 中定义的映射。按该
映射的定义，$\Psi(\sigma_*\gamma)$ 是把 $\tilde x$ 送到
$(\sigma_*\gamma)\curvearrowright\tilde x$ 的万有覆盖 $p:\tilde X\to X$ 的唯一自同构。
但 $(\sigma_*\gamma)\curvearrowright\tilde x=(\widetilde{\sigma^{-1}\circ\gamma})(1)$，
其中 $\widetilde{\sigma^{-1}\circ\gamma}$ 是 $\tilde X$ 中满足
$p\circ\widetilde{\sigma^{-1}\circ\gamma}=\sigma^{-1}\circ\gamma$ 且
$\widetilde{\sigma^{-1}\circ\gamma}(0)=\tilde x$ 的唯一道路。照例用 $\tilde\gamma$ 记
$\tilde X$ 中满足 $p\circ\tilde\gamma=\gamma$ 且 $\tilde\gamma(0)=\tilde x$ 的唯一道路。
由于按 $\tilde\sigma$ 的定义（见引理 2.2.11），
\[
(\tilde\sigma^{-1}\circ\gamma)(0)=\tilde\sigma^{-1}(\tilde\gamma(0))
=\tilde\sigma^{-1}(\tilde x)=\tilde x
\]
且
\[
p\circ(\tilde\sigma^{-1}\circ\tilde\gamma)
=(p\circ\tilde\sigma^{-1})\circ\tilde\gamma
=(\sigma^{-1}\circ p)\circ\tilde\gamma
=\sigma^{-1}\circ(p\circ\tilde\gamma)
=\sigma^{-1}\circ\gamma,
\]
由这种提升的唯一性得 $\widetilde{\sigma^{-1}\circ\gamma}
=\tilde\sigma^{-1}\circ\tilde\gamma$。于是
\[
(\sigma_*\gamma)\curvearrowright\tilde x
=(\tilde\sigma^{-1}\circ\tilde\gamma)(1)
=\tilde\sigma^{-1}(\tilde\gamma(1))
=\tilde\sigma^{-1}(\gamma\curvearrowright\tilde x)
=\tilde\sigma^{-1}(f(\tilde x)),
\]
其中 $f=\Psi(\gamma)$。由于 $\tilde\sigma(\tilde x)=\tilde x$，
$(\sigma_*\gamma)\curvearrowright\tilde x
=(\tilde\sigma^{-1}\circ f)(\tilde x)
=(\tilde\sigma^{-1}\circ f\circ\tilde\sigma)(\tilde x)$。此外，
$\tilde\sigma^{-1}\circ f\circ\tilde\sigma$ 是覆盖 $\tilde X/X$ 的自同构，故由映射
$\Psi$ 的定义 $\Psi(\sigma_*\gamma)=\tilde\sigma^{-1}\circ f\circ\tilde\sigma$。然后
我们注意到，$\Psi(\sigma_*\gamma)$ 实际上等于
$\pi^R_1((X,\sigma);x)\subset\Aut^{\pm}(\tilde X)^{\mathrm{op}}$ 中的乘积
$\tilde\sigma f\tilde\sigma^{-1}$。
\end{proof}

\begin{remark}\label{rem:2.2.15}
注意：若取 $\tilde X$ 为从 $x\in X^\sigma$ 出发的道路的同伦类之通常集合，则自然取
$\tilde x$ 为 $x$ 处常道路的同伦类。特别地，$\tilde\sigma$ 完全由 $x$ 的选取决定，
故命题 2.2.13 与推论 2.2.14 中的半直积典范同构。
\end{remark}

当 $x\notin\Fix(\sigma)$ 时，命题 2.2.17 给出 $\pi^R_1((X,\sigma);x)$ 用 $X$ 中道路的
另一种描述。这需要下述本身也很有意思的引理。

\begin{lemma}\label{lem:2.2.16}
设 $x$ 是 $X$ 中满足 $\sigma(x)\ne x$ 的点，$\eta:[0,1]\to X$ 是从 $x$ 到
$\sigma(x)$ 的道路。由引理 2.2.1，存在唯一的道路 $\tilde\eta:[0,1]\to\tilde X$
使得 $p\circ\tilde\eta=\eta$ 且 $\tilde\eta(0)=\tilde x$。则存在唯一的反全纯映射
$\tilde\sigma:\tilde X\to\tilde X$，使得 $p\circ\tilde\sigma=\sigma\circ p$ 且
$\tilde\sigma(\tilde x)=\tilde\eta(1)$。
\end{lemma}

\begin{proof}
这是引理 2.2.11(i) 的直接应用。注意一般地 $\tilde\sigma^2\ne\Id_{\tilde X}$（相关
讨论另见习题 2.2.30）。
\end{proof}
